\section{Conclusion}
\label{sec:conclusion}

We gave a precise formulation, over a stratified fragment $\CICm$
of the Calculus of Inductive Constructions, of the three
optimisations that separate the current CoqHammer translation from
the state of the art of hammers for simple type theory: heuristic
monomorphisation, bridge-free signature specialisation, and
monotonicity-based guard erasure.  All three were proved sound with
respect to a set-theoretic standard model, via a self-contained
truth lemma for the guarded translation in which the guard
predicate is membership.  The main technical results are: soundness
of monomorphisation by semantic instantiation, independent of the
instance-selection heuristic, with conservativity for the keep
policy and the uncomputability of complete instance selection;
faithfulness of total specialisation on the first-order core
without bridging axioms; and the guard-erasure theorem --- guards
of witnessed, monotonic types may be erased, with monotonicity
certified either by naked-variable analysis or by an internal
infinity criterion that the translation's own injectivity and
discrimination axioms establish for recursive inductive types.
Both side conditions were shown necessary, and the known
counterexample to guard omission in the CoqHammer literature was
identified as a non-monotonicity phenomenon.

Several directions remain open.  \emph{Theory:} the per-variable
(featherweight) refinement of erasure in the single-domain setting;
a polymorphic monotonicity calculus applicable to kept generic
premises; faithfulness for shape-specialised guards over open
types; extending $\CICm$ with indexed inductive families,
cumulativity and fixpoints (we expect the semantic method to scale,
at the cost of a heavier model section); and a mechanisation, for
which the Isabelle formalisation of the monomorphic sorted
metatheory \cite{BlanchettePopescu2013} is a natural starting
point.  \emph{Practice:} implementing the pipeline in CoqHammer
behind options, with the ablation evaluation of
\cref{rem:implementation}; the literature
\cite{Blanchette2016,Desharnais2022} predicts the largest gains
from the combination of drop-policy monomorphisation with erasure,
which is precisely the configuration our theorems license.
